%%%%%%%%%%%%%%%%%%%%%%%%%%%%%%%%%%%%%%%%%%%%%%%%%%%%%%%%%%%%%%%%%%%%%%
% Lớp 11 | Đường thẳng, mặt phẳng. Quan hệ song song trong | Hai đường thẳng song song
% Dạng 1H4?2-6: Ba điểm thẳng hàng
% Số câu: 4  (sinh tự động bởi tach_cau.py; từ nay thầy sửa trực tiếp file này)
%%%%%%%%%%%%%%%%%%%%%%%%%%%%%%%%%%%%%%%%%%%%%%%%%%%%%%%%%%%%%%%%%%%%%%

%%%%%%%%%%%%%%%% ĐÚNG/SAI %%%%%%%%%%%%%%%%

% nguồn: Ok_dungsai_lop11_Lan2
\begin{ex}%[1H4V2-6]%[ma:f8823c47]
	Cho hình chóp $S.ABCD$ có đáy $ABCD$ là hình thang với $AD$ là đáy lớn. Gọi $E$, $F$ lần lượt là trung điểm của $SA$ và $SD$. $K$ là giao điểm của các đường thẳng $AB$ và  $CD$.
	 \choiceTF
		{\True Giao điểm $M$ của đường thẳng  $SB$ và mặt phẳng $(CDE)$ là điểm thuộc đường thẳng $KE$}
		{Đường thẳng $SC$ cắt mặt phẳng $(EFM)$ tại $N$. Tứ  giác $EFNM$ là hình bình hành}
		{\True Các đường thẳng $AM$, $DN$, $SK$ cùng đi qua  một điểm}
		{Cho biết $AD = 2BC$. Tỉ số diện tích của hai tam giác $KMN$ và $KEF$  bằng $\dfrac{2}{3}$}
	\loigiai{ 
		\begin{center}
			\begin{tikzpicture}[scale=0.8,font=\footnotesize,line join=round,line cap=round,>=stealth]
					 \def\a{4}
					 \def\b{1/3}
					\coordinate (O) at (0,0); 
					\coordinate (C) at (\a*1.6,0); 
					\coordinate (A) at  ($(O)+(\a*0.55,\a*0.25)$);
					\coordinate (D) at ($(A)+(C)-(O)$); 
					\coordinate (S) at ($(\a*1.2,\a*1.8)$); 
					\coordinate (B)  at ($(O)!0.5!(C)$);
					 \coordinate (K) at ($(B)!-1!(A)$);
					\coordinate (E) at  ($(S)!0.5!(A)$);
					\coordinate (F) at  ($(S)!0.5!(D)$);
					\coordinate  (M) at ($(E)!\b!(K)$);
					\coordinate (N) at ($(F)!\b!(K)$); 
					\coordinate (I)  at ($(A)!1.5!(M)$);
					\draw (S)--(K)--(D)--(S)--(A)--(K)--(E) 
						(B)--(S)--(C) 
						(A)--(I)--(D); 
					\draw[dashed]  (A)--(D)
						 (F)--(B)--(C)--(E)--(F)--(M)--(N);
					\foreach \x/\label/\g  in {A/\(A\)/-110,B/\(B\)/-110,C/\(C\)/-90,D/\(D\)/-60,S/\(S\)/90,K/\(K\)/-100,I/\(I\)/30,M/\(M\)/180,N/\(N\)/30,E/\(E\)/180,F/\(F\)/30}
					\fill (\x) circle (1pt) 
						 ($(\x)+(\g:3mm)$) node {\small \label};
				\end{tikzpicture}
		\end{center}
		\begin{itemchoice}
			\itemch \textbf{Đúng}  \\
			Ta  có $SK = (SAB) \cap (SCD)$. \\
			Trong mặt phẳng $(SAB)$\\ 
			 \[
			\heva{&M = KE \cap SB \\&KE \subset (CDE)} 
			\Rightarrow SB \cap (CDE) = M. 
			 \]
			\itemch \textbf{Sai}  \\
			 Trong mặt phẳng $(SCD)$\\
			\[ 
			\heva{ 
			&N = KF  \cap SC \\
			&KF  \subset (EFM)
			} 
			 \Rightarrow SC \cap (EFM) = N.
			\] 
			 Ta có
			 \[
			 \heva{
			&MN  = (EFK) \cap (SBC) \\
			&EF  \parallel BC\\
			&EF \subset  (EFK)\\
			&BC \subset (SBC)} 
			 \Rightarrow MN \parallel EF \parallel BC.
			 \]
			Suy ra tứ  giác $EFNM$ là hình thang.
			\itemch  \textbf{Sai}\\
			Trong mặt phẳng $(ADNM)$, gọi $I = AM \cap DN$.  \\
			Ta  có
			 \[
			\heva{&I \in AM, AM \subset (SAB)\\ &I \in CD, CD \subset  (SCD)}
			 \Rightarrow I \in (SAB) \cap (SCD)
			\Rightarrow I \in SK. 
			\] 
			Vậy $3$ đường thẳng $AM, DN, SK$ đồng quy tại điểm $I$. 
			\itemch \textbf{Sai} \\ 
			Khi $AD = 2BC$, dễ dàng chứng minh được $B$, $C$ lần  lượt là trung điểm của $KA$ và $KD$. \\
			Suy ra $M$, $N$  lần lượt là trọng tâm của hai tam giác $SAK$ và $SDK$.\\
			Do đó $MN =  \dfrac{2}{3}EF$.\\
			Gọi $h_1$, $h_2$ lần lượt  là khoảng cách từ điểm $K$ đến $MN$ và $EF$.\\
			Khi đó $h_1 =  \dfrac{2}{3}h_2$.\\
			Vậy 
			\[ 
			\dfrac{S_{KMN}}{S_{KEF}} = \dfrac{\dfrac{1}{2}MN \cdot h_1}{\dfrac{1}{2}EF \cdot h_2} =  \dfrac{\dfrac{2}{3}EF \cdot \dfrac{2}{3}h_2}{EF \cdot h_2} = \dfrac{4}{9}.
			 \]
		\end{itemchoice}
	} 
\end{ex}


%%%%%%%%%%%%%%%% TRẢ LỜI NGẮN %%%%%%%%%%%%%%%%

% nguồn: Ok_TraLoiNgan_Lop11_Lan1
\begin{ex}%[1H4V2-6]%[ma:f5f7d9bb]
	Cho hình chóp $S.ABCD$ có đáy là hình bình hành, cạnh $AB=6$ cm, $AC$ và $BD$ cắt nhau tại $O$. Gọi $I$ là trung điểm của $SO$. Mặt phẳng $(ICD)$ cắt $SA$, $SB$ lần lượt tại $M$, $N$. Độ dài  $MN$ là
	\begin{center}
		\begin{tikzpicture}[scale=0.5,font=\footnotesize,line join=round,line cap=round,>=stealth]
				 \path
					 (0,0) coordinate (A)
					++(-140:3) coordinate  (B)
					++(0:8) coordinate  (C)
					 ($(A)+(C)-(B)$) coordinate (D)
					 ($(A)+(1,6)$) coordinate (S)
					($(A)!0.5!(C)$)  coordinate (O)
					($(S)!0.5!(O)$)  coordinate (I);
				\foreach \i in{B,C,D}{\draw  (S)--(\i);};
				\draw (B)--(C)--(D); 
				\draw[dashed]  (O)--(S)--(A)--(B)--(D)--(A)--(C)
					(D)--(I)--(C); 
				\foreach \i/\g  in {A/-90,B/-90,C/-90,D/-90,S/90,O/-90,I/30}
				\fill[black] (\i) circle  (1pt)+(\g:4mm) node[scale=1]{$\i$};
			\end{tikzpicture}
	\end{center}
	 \shortans{2}
	\loigiai{ 
		\begin{center}
			\begin{tikzpicture}[scale=0.5,font=\footnotesize,line join=round,line cap=round,>=stealth]
					\path 
						 (0,0) coordinate (A)
						++(-140:3)  coordinate (B)
						++(0:8) coordinate  (C)
						($(A)+(C)-(B)$) coordinate (D) 
						 ($(A)+(1,6)$) coordinate (S)
						($(A)!0.5!(C)$) coordinate (O) 
						($(S)!0.5!(O)$) coordinate (I) 
						 ($(S)!1/3!(A)$) coordinate (M)
						($(S)!1/3!(B)$)  coordinate (N);
					\foreach \i in{B,C,D}{\draw (S)--(\i);}; 
					\draw (B)--(C)--(D); 
					 \draw[dashed] (O)--(S)--(A)--(B)--(D)--(A)--(C)
						(D)--(I)--(C) 
						 (I)--(M)--(N)--(I);
					\foreach \i/\g  in {A/-90,B/-90,C/-90,D/-90,S/90,O/-90,I/30,N/150}
					\fill[black] (\i)  circle (1pt)+(\g:4mm) node[scale=1]{$\i$};
					\fill[black] (M) circle (1pt)+(0:4mm)  node[scale=1]{$M$};
				\end{tikzpicture}
		\end{center}
		Ta có $\heva{&I \in  (ICD)\\&I \in SO \subset (SAC)\\&C \in (ICD)\\&C \in (SAC)} \Rightarrow IC$ là giao tuyến của $(ICD)$ và $(SAC)$.\\
		Mặt khác điểm $M$ là giao điểm của $SA$ và $(ICD)$, mà $SA \subset (SAC)$ nên $M$ nằm  trên giao tuyến của $(ICD)$ và $(SAC)$.\\
		Suy ra $M$,  $I$, $C$ thẳng hàng.
		Tương tự,  ta cũng chứng minh được $N$, $I$, $D$ thẳng hàng.\\
		Áp dụng Định lý Menelaus  trong $\triangle SAO$\\
		$$\dfrac{CA}{CO}\cdot  \dfrac{IO}{IS} \cdot \dfrac{MS}{MA} = 1\Rightarrow 2\cdot 1\cdot \dfrac{MS}{MA} = 1\Rightarrow \dfrac{MS}{MA} = \dfrac{1}{2}\Rightarrow \dfrac{SM}{SA} = \dfrac{1}{3}.$$
		Ta thấy $AB$, $CD$, $MN$ lần lượt là giao  tuyến của ba mặt phẳng $(SAB)$, $(SCD)$ và $(ICD)$.\\
		Theo định lý về giao tuyến của ba mặt phẳng và $AB \parallel CD$ nên $MN \parallel AB \parallel CD$.\\ 
		Áp dụng Định lý Thales  trong $\triangle SAB$
		$$\dfrac{MN}{AB} = \dfrac{SM}{SA} = \dfrac{1}{3}\Rightarrow MN = \dfrac{1}{3}\cdot  AB = \dfrac{1}{3}\cdot 6 = 2.$$
	 }
\end{ex}

% nguồn: Ok_TraLoiNgan_Lop11_Lan2
\begin{ex}%[1H4V2-6]%[ma:65eb9893]
	Cho tứ diện $ABCD$ có các cạnh bằng nhau và bằng $10$. Gọi $E$ là điểm thuộc cạnh $AB$ sao cho $AE=2EB$, $F$ là điểm thuộc cạnh $BC$ sao cho $BF=2FC$ và $G$ là điểm thuộc cạnh $CD$ sao cho $CG=2GD$. Tính độ dài đoạn giao tuyến của mặt phẳng $(EFG)$ và mặt bên $ACD$ (làm tròn kết quả đến hàng phần trăm). 
	 \shortans{1,47}
	\loigiai{ 
		\begin{center}
			\begin{tikzpicture}[scale=1,font=\footnotesize,line join=round,line cap=round,>=stealth]
					 \coordinate[label=left:$B$] (B) at (0,0);
					\coordinate[label=right:$D$] (D) at (4,0); 
					\coordinate[label=below  left:$C$] (C) at (-50:3);
					\coordinate[label=above:$A$] (A)  at (70:4);
					\coordinate[label=left:$E$] (E)  at ($(A)!2/3!(B)$);
					\coordinate[label=left:$F$] (F) at  ($(B)!2/3!(C)$);
					\coordinate[label=right:$M$] (M)  at (5,0);
					\coordinate[label=above right:$N$]  (N) at (intersection of E--M and A--D);
					 \coordinate[label=below:$G$] (G) at (intersection of F--M and C--D);
					\draw (B)--(C)--(G)--(N)--(A)--cycle  (A)--(C)
						 (E)--(F)
						(M)--(G)--(N)--cycle; 
					\draw[dashed]  (B)--(D)
						 (E)--(N)--(D)--(G)--(F);
					 \foreach \diem in {B,C,D,A,E,F,M,N,G} \fill (\diem) circle (1pt);
				\end{tikzpicture}
		\end{center}
		Gọi $M=FG\cap BD$  và $N=EM\cap AD$.\\
		Suy  ra $NG=(EFG)\cap(ACD)$.\\
		Áp dụng định lý Menelaus cho  tam giác $BDC$ với ba điểm $F$, $G$, $M$ thẳng hàng, ta có
		 \[\dfrac{MD}{MB}\cdot\dfrac{FB}{FC}\cdot\dfrac{GC}{GD}=1\Rightarrow \dfrac{MD}{MB}\cdot2\cdot2=1\Rightarrow \dfrac{MD}{MB}=\dfrac{1}{4} \quad(1)\]
		Từ $(1)$ và áp dụng định lý Menelaus cho tam giác $ABD$ với ba điểm $E$, $M$, $N$ thẳng hàng,  ta được
		 \[\dfrac{MB}{MD}\cdot\dfrac{ND}{NA}\cdot\dfrac{EA}{EB}=1\Rightarrow 4\cdot\dfrac{ND}{NA}\cdot2=1\Rightarrow \dfrac{ND}{NA}=\dfrac{1}{8}\quad (2)\]
		 Từ $(2)$ ta được
		$8ND=NA\Rightarrow 9ND=10\Rightarrow ND=\dfrac{10}{9}$.\\ 
		Xét tam  giác đều $NDG$, ta có
		\[NG=\sqrt{ND^2+DG^2-2ND\cdot DG\cdot\cos60^{\circ}}=\sqrt{\left(\dfrac{10}{9}\right)^2+\left(\dfrac{10}{3}\right)^2-2\cdot\dfrac{10}{9}\cdot\dfrac{10}{3}\cdot\dfrac{1}{2}}=\dfrac{5\sqrt{7}}{9}\approx 1{,}47.\] 
	} 
\end{ex}

% nguồn: Ok_TraLoiNgan_Lop11_Lan2
\begin{ex}%[1H4V2-6]%[ma:25ba4cc9]
	Cho tứ diện $ABCD$ có các cạnh bằng nhau và bằng $ 10$. Gọi $E$ là điểm thuộc cạnh $AB$ sao cho $AE=2EB$, $ F$ là điểm thuộc cạnh $BC$ sao cho $BF=2FC$ và $ G$ là điểm thuộc cạnh $CD$ sao cho $CG=2GD$. Độ dài đoạn giao tuyến của mặt  phẳng $\left(EFG\right)$ và mặt bên $(ACD)$ bằng
	\shortans{1,47} 
	\loigiai{ 
		\begin{center}
			\begin{tikzpicture}[scale=1,line join=round,line cap=round,font=\footnotesize,>=stealth]
					\def\a{4} 
					\path (0:0) coordinate  (B)
						++(0:\a) coordinate  (D)
						++(-120:\a/2)  coordinate (C)
						($(B)+(70:\a)$) coordinate  (A)
						($(A)!2/3!(B)$) coordinate (E) 
						($(B)!2/3!(C)$) coordinate  (F)
						($(C)!2/3!(D)$) coordinate  (G)
						 ($(A)!8/9!(D)$) coordinate (N)
						(intersection of E--N and F--G) coordinate  (M);
					 \draw[dashed,thick] (B)--(D)--(M)
						 (E)--(N)
						 (F)--(G)--(D)--(N);
					 \draw[thick] (B)--(A)--(N)
						(A)--(C) 
						(B)--(C)--(G)--(M)--(N) 
						 (E)--(F);
					\foreach  \x/\g in {A/90,B/180,C/-45,D/-20,E/20,F/-90,G/-45,N/45,M/0}
					\fill[black] (\x) circle (1pt) 
						($(\g:3mm)+(\x)$)  node {$\x$};
				\end{tikzpicture}
		\end{center}
		Gọi $\left\{ M\right\}=FG\cap BD$ và $\left\{ N\right\}=EM\cap AD$. Suy  ra $ NG=\left(EFG\right)\cap\left(ACD\right)$.\\
		Theo định lí  Menelaus cho tam giác $BDC$ với 3 điểm $F$, $G$, $M$ thẳng hàng ta có
		$$\dfrac{MD}{MB} \cdot \dfrac{FB}{FC} \cdot  \dfrac{GC}{GD}= 1 \Rightarrow \dfrac{MD}{MB}\cdot 2\cdot 2=1 \Rightarrow\dfrac{MD}{MB}=\dfrac{1}{4}.$$
		Tương tự, ta có $\dfrac{MB}{MD} \cdot \dfrac{ND}{NA}\cdot \dfrac{EA}{EB}=1 \Rightarrow 4\cdot 2 \cdot  \dfrac{ND}{NA}=1 \Rightarrow\dfrac{ND}{NA}=\dfrac{1}{8}$.\\
		$\Rightarrow  8ND=NA\Rightarrow 9ND=10\Rightarrow ND=\dfrac{10}{9}$.\\
		Xét tam  giác đều $ NDG$, ta có\\
	 $$ NG=\sqrt{N{D^2}+D{G^2}-2ND \cdot DG \cdot \cos 60^\circ}=\sqrt{\dfrac{10^2}{81}+\dfrac{10^2}{9}-2 \cdot \dfrac{10}{9}\cdot \dfrac{10}{3}\cdot \dfrac{1}{2}}=\dfrac{5\sqrt 7}{9}\approx 1{,}47.$$}
\end{ex}
